\chapter{研究内容}

\section{研究目标}

本项目面向动态未知环境与非理想通信条件下无人机自主飞行需求，围绕动态障碍预测可靠性和多无人机协同信息可靠性开展研究，形成“动态障碍短时预测与不确定性表征—风险感知运动基元局部规划—通信受限多无人机自适应协同管理”的技术体系，提升无人机在复杂动态环境中的安全性、实时性以及通信受限条件下的协同鲁棒性。

具体目标包括：

（1）建立动态障碍短时运动预测及预测不确定性表征方法，使规划器不仅获得障碍未来状态预测结果，同时能够评价预测结果的可信范围。

（2）建立面向预测不确定性的轻量化风险感知运动基元规划方法，使预测不确定性通过低复杂度时空风险直接参与候选轨迹筛选，同时满足高频在线规划的实时性要求。

（3）建立面向通信受限环境的邻机相关性感知与协同约束自适应机制，使多无人机能够根据通信状态和邻机冲突相关性动态筛选协同对象，并自适应调整轨迹可信时间窗、安全裕度和冲突约束。

\section{主要研究内容}

本研究围绕“环境状态预测—单机局部安全规划—多机通信约束协同”三个层次展开。第一部分解决动态障碍未来状态及其可信程度如何获得的问题；第二部分研究预测不确定性如何以较低计算代价进入运动基元筛选和安全决策；第三部分在单机局部规划基础上扩展至多无人机场景，研究通信状态、邻机协同相关性与规划约束之间的动态关联。

\subsection{动态障碍运动预测及不确定性表征}

针对动态未知环境中障碍未来运动状态难以准确获得的问题，研究基于连续局部感知信息的动态障碍状态估计和短时运动预测方法。利用LiDAR等机载传感器获得周围环境观测，对动态障碍进行目标提取、数据关联和状态跟踪，形成障碍物位置、速度等状态量。令第$k$时刻动态障碍状态为

\begin{equation}
\mathbf{x}_{o,k}=\begin{bmatrix}\mathbf{p}_{o,k}^{\mathrm T} & \mathbf{v}_{o,k}^{\mathrm T}\end{bmatrix}^{\mathrm T},
\end{equation}

其中，$\mathbf{p}_{o,k}$和$\mathbf{v}_{o,k}$分别表示障碍物位置和速度。基于短时运动模型获得预测时域内的未来状态：

\begin{equation}
\hat{\mathbf{x}}_{o,k+\tau}=f(\mathbf{x}_{o,k},\tau),
\end{equation}

其中，$\tau$表示预测时间间隔，$f(\cdot)$表示短时状态预测模型。

对于人员、移动机器人等具有自主运动能力的动态障碍，其短时运动可能出现转向、加减速或运动模式突变；同时局部传感器还会受到噪声、遮挡和有限历史观测影响，因此由匀速或局部线性模型获得的单一未来轨迹不可避免存在误差。若规划器直接将该预测结果视为确定状态进行碰撞判断，容易在预测失准时低估碰撞风险。因此，本研究进一步采用预测均值与协方差或等价安全包络描述未来状态：

\begin{equation}
\mathbf{x}_{o,k+\tau}\sim\mathcal{N}\left(\boldsymbol{\mu}_{o,\tau},\boldsymbol{\Sigma}_{o,\tau}\right).
\end{equation}

重点研究预测误差随预测时间增长的传播规律，并结合障碍物几何尺寸，将预测均值和不确定性转化为能够被局部规划器快速查询的未来危险区域或时空风险表示。该部分主要输出“未来障碍状态+预测不确定性”，作为后续风险感知运动基元规划的统一输入。SOGM未来时空占据表示、碰撞概率和风险评估等研究为该部分提供方法基础\cite{thomas2022sogm,thomas2023future,liu2023tight,xia2024risk}。

\subsection{面向预测不确定性的轻量化风险感知运动基元规划}

现有动态局部规划可概括为两类思路。一类依据当前观测或单一未来预测轨迹进行快速碰撞检测，计算代价低，但容易对未来状态产生过度确定的假设；另一类将协方差、碰撞概率、机会约束或风险包络显式纳入规划，能够描述预测误差，但在连续优化或复杂概率计算条件下可能增加在线计算负担。基于此，本研究拟在运动基元框架中建立低复杂度不确定性风险评价，使风险真正改变候选轨迹选择，而不直接采用高开销的全局概率优化。

根据无人机当前状态、速度范围和动力学约束生成候选运动基元集合：

\begin{equation}
\mathcal{M}=\{M_1,M_2,\ldots,M_N\}.
\end{equation}

对于第$j$条候选运动基元，建立综合评价函数：

\begin{equation}
J_j=w_gJ_g+w_sJ_s+w_eJ_e+w_rJ_r,
\end{equation}

其中，$J_g$表示目标趋近代价，$J_s$表示轨迹平滑代价，$J_e$表示运动或能量代价，$J_r$表示动态环境风险代价，$w_g,w_s,w_e,w_r$为相应权重。

风险项综合考虑候选轨迹与障碍预测均值之间的空间距离、时间关系及预测不确定程度：

\begin{equation}
J_r=\sum_{\tau=1}^{H}\Phi\left(\mathbf{p}_{j,\tau},\boldsymbol{\mu}_{o,\tau},\boldsymbol{\Sigma}_{o,\tau}\right),
\end{equation}

其中，$H$表示预测时域，$\mathbf{p}_{j,\tau}$表示候选轨迹在未来时刻的空间位置。当候选轨迹接近障碍未来高风险区域，或障碍预测不确定程度增大时，提高其风险代价或扩大对应风险边界，使规划器主动选择具有更高安全裕度的运动基元。

为避免单纯依赖软代价在极端状态下失效，进一步引入碰撞时间（Time-to-Collision，TTC）、最小安全距离（clearance）等硬安全约束。对满足硬安全条件的候选轨迹进行风险排序；当正常候选整体不满足安全要求时，进入减速、制动、侧移或爬升等应急动作集合，形成“低复杂度风险软筛选+TTC/安全距离硬过滤”的双层安全机制。运动基元实时规划、机会约束运动基元及碰撞预测等研究为该方案提供基础\cite{primitiveSwarm,gutow2020koopman,nguyen2022collision}。

\subsection{通信受限下邻机相关性感知与协同约束自适应规划}

针对非理想通信条件下邻机共享状态滞后、轨迹信息失效以及不同邻机对当前规划作用不同的问题，研究通信状态、邻机相关性和协同约束之间的联合调节机制。对于无人机$i$从邻机$j$接收到的信息，首先定义通信状态向量：

\begin{equation}
\mathbf{s}^{\mathrm{comm}}_{ij}=\begin{bmatrix}\tau_{ij} & p^{\mathrm{loss}}_{ij} & \Delta_{ij} & q_{ij}\end{bmatrix}^{\mathrm T},
\end{equation}

其中，$\tau_{ij}$表示通信时延，$p^{\mathrm{loss}}_{ij}$表示数据丢失特征，$\Delta_{ij}$表示距最近一次有效信息更新的时间，$q_{ij}$表示链路质量。信息年龄（Age of Information，AoI）可作为信息新鲜程度的重要量化基础，但本研究不将AoI或通信质量本身直接等价为协同权重，而是将其作为协同管理的输入之一\cite{hu2024aoi}。

进一步结合规划层因素构造邻机协同相关性：

\begin{equation}
R_{ij}=G\left(\mathbf{s}^{\mathrm{comm}}_{ij},d_{ij},TTC_{ij},\chi_{ij}\right),\qquad R_{ij}\in[0,1],
\end{equation}

其中，$d_{ij}$表示机间距离，$TTC_{ij}$表示基于相对运动估计的潜在碰撞时间，$\chi_{ij}$表示轨迹交叉或冲突趋势特征。该相关性用于区分“通信状态差但与当前规划无关的邻机”和“通信信息开始过期但具有较高冲突风险的关键邻机”，从而避免所有邻机采用统一处理策略。

根据通信状态与邻机相关性联合筛选需要重点参与协同冲突判断的邻机集合，并动态调整邻机轨迹可信时间窗、机间安全裕度、状态外推膨胀范围和冲突约束强度。当通信状态良好且关键邻机信息保持新鲜时维持正常协同；当信息开始过期但仍具有规划价值时进入保守协同；当关键信息严重失效时降低对共享轨迹的直接依赖，并切换至以本机局部感知、邻机短时外推和硬安全约束为主的局部自主安全模式。

最终形成“通信状态评价—邻机相关性筛选—协同约束自适应—必要时模式退化—局部安全规划”的闭环机制。该研究点与已有通信对象选择、通信感知路径规划和时延鲁棒规划形成联系，但研究重点放在通信退化后哪些邻机仍值得协同，以及如何调整规划约束，而不是单独优化通信指标\cite{serragomez2020whom,robustMader,guo2024communication,liu2024power,xu2025search}。

\section{拟解决的关键问题}

\subsection{动态障碍预测误差如何在线表征并转化为可用于碰撞判断的风险边界}

动态障碍未来状态受观测误差、模型偏差和运动行为变化共同影响，需要研究如何利用有限历史观测估计未来状态及可信范围，并将预测均值、协方差或安全包络转化为规划器可实时查询的风险边界，避免单一确定性预测在运动突变时造成碰撞风险低估。

\subsection{如何在满足实时性的条件下使预测不确定性直接影响运动基元选择}

复杂概率规划能够提高不确定环境中的安全性，但计算负担较高。本研究需要解决如何以低复杂度方式把预测不确定性嵌入运动基元风险评价，使风险真正改变候选轨迹排序与筛选，同时维持高频在线规划，并通过硬安全层对软风险评价进行兜底。

\subsection{通信受限条件下如何筛选有效邻机并自适应调整多机协同约束}

通信时延、丢包和信息更新时间属于通信层指标，但多机规划实际需要判断的是哪些邻机信息仍具有决策价值、哪些邻机当前存在较高冲突风险，以及通信退化后应如何调整安全约束。需要建立通信状态、邻机相关性、协同对象筛选与规划约束调整之间的映射，使通信质量变化能够直接改变多机局部规划行为，而不是仅作为独立网络性能指标。